\documentclass[11pt,a4paper]{article}

\usepackage[margin=1.1in]{geometry}
\usepackage[T1]{fontenc}
\usepackage[utf8]{inputenc}
\usepackage{lmodern}
\usepackage{amsmath,amssymb,amsthm,mathtools,mathrsfs}
\usepackage{enumitem}
\usepackage{microtype}
\usepackage[colorlinks=true,linkcolor=blue,citecolor=blue,urlcolor=blue]{hyperref}

\newtheorem{theorem}{Theorem}[section]
\newtheorem{proposition}[theorem]{Proposition}
\newtheorem{lemma}[theorem]{Lemma}
\newtheorem{corollary}[theorem]{Corollary}
\theoremstyle{definition}
\newtheorem{definition}[theorem]{Definition}
\theoremstyle{remark}
\newtheorem{remark}[theorem]{Remark}

\DeclareMathOperator{\Fix}{Fix}
\DeclareMathOperator{\Ann}{Ann}
\DeclareMathOperator{\Spec}{Spec}
\DeclareMathOperator{\ord}{ord}
\DeclareMathOperator{\Idl}{Idl}
\DeclareMathOperator{\Tail}{Tail}
\DeclareMathOperator{\supp}{supp}
\DeclareMathOperator{\im}{im}
\DeclareMathOperator{\Arg}{Arg}
\DeclareMathOperator{\width}{width}

\newcommand{\ZZ}{\mathbb Z}
\newcommand{\NN}{\mathbb N}
\newcommand{\RR}{\mathbb R}
\newcommand{\CC}{\mathbb C}
\newcommand{\PP}{\mathbb P}
\newcommand{\BB}{\mathbf B}
\newcommand{\QQ}{\mathbb Q}
\newcommand{\cO}{\mathcal O}
\newcommand{\cP}{\mathcal P}
\newcommand{\cV}{\mathcal V}
\newcommand{\cA}{\mathcal A}
\newcommand{\cB}{\mathcal B}
\newcommand{\cQ}{\mathcal Q}
\newcommand{\cK}{\mathcal K}
\newcommand{\cI}{\mathcal I}
\newcommand{\cC}{\mathcal C}
\newcommand{\cE}{\mathcal E}
\newcommand{\taup}{\tau}
\newcommand{\into}{\hookrightarrow}
\newcommand{\onto}{\twoheadrightarrow}

\title{Sheafified Split-Zero Shadows of the Zeta Zero Divisor\\
Boolean Cubes, Jet-Ideal Chains, Compactified Tails, and Packet Descent}
\author{}
\date{May 2026}

\begin{document}
\maketitle

\begin{abstract}
Let
\[
S=G(\ZZ)=\ZZ\sqcup\{\taup\},\qquad A:=\{0,\taup\}\cong\BB,
\qquad F_1(s):=\zeta(s)-1.
\]
For every zero $\rho$ of the Riemann zeta function one has $F_1(\rho)=-1$, and multiplication by $-1$ fixes exactly the Boolean pair $A$.  The pointwise split-zero datum is therefore constant.  The purpose of the paper is to show that the zero divisor nonetheless carries a nontrivial collection of sheaf-theoretic and infinitesimal invariants.

For an effective discrete divisor $D=\sum_x m_x[x]$ on a Riemann surface, we prove that the finite-support presheaf of local split-zero data sheafifies to the Boolean multiplicity shadow
\[
\cV_D(U)=\prod_{x\in U\cap |D|}A^{m_x},
\]
with reduced retract
\[
\cA_D(U)=\prod_{x\in U\cap |D|}A.
\]
We identify the ideal lattice of the classical local Artin thickening $\cO_{X,x}/\mathfrak m_x^{m_x}$ with a finite chain $C_{m_x}$, embed that chain into the Boolean cube $B_{m_x}=\cP([m_x])$, and prove the explicit biadjoint relation
\[
L_{m_x}\dashv i_{m_x}\dashv R_{m_x}.
\]
For a holomorphic function $f$ with divisor $D$, the local class of $f$ in the $(m_x+1)$-jet algebra generates the minimal nonzero ideal; on $\CC$ this recovers the coefficient $f^{(m_x)}(x)/m_x!$.  After compactification to $\PP^1(\CC)$ we compute the defect between direct image and extension by zero and identify the tail object at $\infty$.  For the completed zeta function $\xi$, the nontrivial-zero shadows descend to packet sheaves on the Klein-four orbit space.  On the trivial-zero sector we recover the exact formula
\[
|\zeta'(-2n)|=\frac{(2n)!}{2(2\pi)^{2n}}\,\zeta(2n+1).
\]
The final section records explicit cyclotomic phase-coordinate lemmas, based on unpublished notes of Jacolm Tobley, for normalized affine factors in the CM field $\QQ(\zeta_5)$ and for their elementary-symmetric expansions.
\end{abstract}

\tableofcontents

\section{Introduction}

Let
\[
S=G(\ZZ)=\ZZ\sqcup\{\taup\},\qquad A:=\{0,\taup\}\subset S,
\qquad F_1(s):=\zeta(s)-1.
\]
If $\rho$ is a zero of $\zeta$, then
\[
F_1(\rho)=\zeta(\rho)-1=-1.
\]
In the globalization semiring, multiplication by $-1$ fixes precisely the Boolean pair $A\cong\BB$.  Pointwise, therefore, every zeta zero determines the same split-zero fiber.  The question is not pointwise identification but global and infinitesimal organization of the zero divisor.

The paper proves seven families of results.

\begin{enumerate}[label=(\arabic*)]
\item The Boolean multiplicity shadow is not inserted ad hoc: it is the sheafification of the finite-support presheaf of local split-zero data.
\item The reduced shadow is a canonical split retract of the multiplicity shadow, and the two coincide if and only if all multiplicities are $1$.
\item The local Artin thickening $\cO_{X,x}/\mathfrak m_x^{m_x}$ has ideal lattice equal to a finite chain $C_{m_x}$, while the Boolean shadow stalk is the cube $B_{m_x}=\cP([m_x])$; the precise comparison is the biadjunction
\[
L_{m_x}\dashv i_{m_x}\dashv R_{m_x}.
\]
\item For a holomorphic function $f$ with divisor $D$, the local jet class in $\cO_{X,x}/\mathfrak m_x^{m_x+1}$ generates the minimal nonzero ideal.  On $\CC$ there is also a coefficient-support retraction from a derivative-enhanced pointed-set sheaf to the Boolean cube.
\item After compactification to $\PP^1(\CC)$, the difference between direct image and extension by zero is measured exactly by a tail object at $\infty$.
\item The nontrivial zeros of the completed zeta function descend, with multiplicities, to packet sheaves on the Klein-four orbit space.
\item The trivial zeros give simple Boolean stalks with exact analytic weight $|\zeta'(-2n)|$, and the final section supplies explicit cyclotomic phase coordinates and orthogonal-walk expansions for normalized local jet phases.
\end{enumerate}

The geometric background is standard: local analytic geometry on Riemann surfaces \cite{Lee}, affine and projective constructions in algebraic geometry \cite{GrothendieckEGAII,Vakil,OrecchiaChiantini}, semirings and idempotent algebra \cite{GaubertKatz,Gunawardena,DrosteKuichVogler}, and the blueprint formalism \cite{Lorscheid}.  The proofs below are self-contained.

\begin{remark}
Section~\ref{sec:tobley} is based on unpublished notes of Jacolm Tobley on cyclotomic affine factors, arctangent decompositions, and orthogonal-walk expansions.  The source transcripts supplied with the project contain orthographic variants of the name; we use the attribution specified with the source bundle.
\end{remark}

\section{Globalization semirings and the fixed Boolean pair}

\begin{definition}
Let $R$ be a commutative ring with identity.  Its \emph{globalization semiring} is
\[
G(R):=R\sqcup\{\taup\},
\]
with operations
\[
a\oplus b:=a+b,\qquad a\oplus\taup=\taup\oplus a:=a,\qquad \taup\oplus\taup:=\taup,
\]
and
\[
a\odot b:=ab,\qquad a\odot\taup=\taup\odot a:=\taup,\qquad \taup\odot\taup:=\taup.
\]
Thus $\taup$ is the additive identity and multiplicative absorber, while $1\in R\subset G(R)$ is the multiplicative identity.
\end{definition}

\begin{definition}
For $u\in R^\times\subset G(R)^\times$, define
\[
\Fix_{G(R)}(u):=\{x\in G(R):u\odot x=x\}.
\]
\end{definition}

\begin{lemma}\label{lem:units-fix}
For every commutative ring $R$,
\[
G(R)^\times=R^\times.
\]
Moreover, for every $u\in R^\times$,
\[
\Fix_{G(R)}(u)=\{\taup\}\cup\Ann_R(u-1).
\]
\end{lemma}

\begin{proof}
If $x\in G(R)$ is invertible, say $x\odot y=1$, then $x\neq\taup$ because $\taup\odot y=\taup\neq1$.  Likewise $y\neq\taup$.  Hence $x,y\in R$ and $xy=1$ inside $R$, so $x\in R^\times$.  The converse is immediate.

Now let $u\in R^\times$.  Since $u\odot\taup=\taup$, the absorber is always fixed.  For $r\in R$ one has
\[
u\odot r=r \iff ur=r \iff (u-1)r=0,
\]
which proves the formula.
\end{proof}

\begin{proposition}\label{prop:split-root-domain}
Let $R$ be an integral domain and let $u\in R^\times$ satisfy $u\neq1$.  Then
\[
\Fix_{G(R)}(u)=\{0,\taup\}.
\]
In particular, if $\xi\in\cO_K^\times$ is a nontrivial root of unity in the ring of integers of a number field $K$, then
\[
\Fix_{G(\cO_K)}(\xi)=\{0,\taup\}.
\]
\end{proposition}

\begin{proof}
By Lemma~\ref{lem:units-fix},
\[
\Fix_{G(R)}(u)=\{\taup\}\cup\Ann_R(u-1).
\]
Because $R$ is a domain and $u-1\neq0$, multiplication by $u-1$ is injective.  Hence $\Ann_R(u-1)=\{0\}$.
\end{proof}

\begin{proposition}\label{prop:boolean-pair}
For every ring $R$, the subset
\[
A_R:=\{0,\taup\}\subset G(R)
\]
is a subsemiring canonically isomorphic to the Boolean semiring $\BB=\{0,1\}$.  Explicitly,
\[
\taup\longleftrightarrow 0\in\BB,
\qquad
0\longleftrightarrow 1\in\BB.
\]
\end{proposition}

\begin{proof}
Inside $A_R$ the operation tables are
\[
0\oplus0=0,\qquad 0\oplus\taup=0,\qquad \taup\oplus\taup=\taup,
\]
and
\[
0\odot0=0,\qquad 0\odot\taup=\taup,\qquad \taup\odot\taup=\taup.
\]
These are exactly the Boolean rules after the displayed identification.
\end{proof}

\begin{corollary}\label{cor:zeta-pointwise}
Let $S=G(\ZZ)$ and $F_1(s)=\zeta(s)-1$.  If $\rho$ is any zero of $\zeta$, then
\[
F_1(\rho)=-1,
\qquad
\Fix_S(F_1(\rho))=\{0,\taup\}\cong\BB.
\]
\end{corollary}

\begin{proof}
The first identity is immediate from $\zeta(\rho)=0$.  The second is Proposition~\ref{prop:split-root-domain} with $R=\ZZ$ and $u=-1$.
\end{proof}

\section{Finite-support presheaves and Boolean divisor shadows}

Throughout this section, $X$ is a connected Riemann surface and
\[
D=\sum_{x\in Z}m_x[x]
\]
is an effective divisor with discrete support $Z\subset X$.  We write $A:=\{0,\taup\}\cong\BB$.

\begin{definition}
The \emph{finite-support presheaf} of split-zero data attached to $D$ is
\[
\cP_D(U):=\bigoplus_{x\in U\cap Z}A^{m_x},
\]
that is, the set of families $(a_x)_{x\in U\cap Z}$ with $a_x\in A^{m_x}$ and all but finitely many $a_x$ equal to the basepoint $(\taup,\dots,\taup)$.  Restriction maps are coordinate projection.

The \emph{Boolean multiplicity shadow} of $D$ is the sheaf of $A$-semimodules
\[
\cV_D(U):=\prod_{x\in U\cap Z}A^{m_x},
\]
and the \emph{reduced Boolean shadow} is
\[
\cA_D(U):=\prod_{x\in U\cap Z}A.
\]
\end{definition}

\begin{lemma}\label{lem:product-sheaf}
Let $Z\subset X$ be closed and discrete, and let $(M_x)_{x\in Z}$ be a family of sets, posets, semimodules, or vector spaces.  Then
\[
\mathcal F_M(U):=\prod_{x\in U\cap Z}M_x
\]
defines a sheaf on $X$.
\end{lemma}

\begin{proof}
Let $U=\bigcup_{i\in I}U_i$ be an open cover.  If two sections of $\mathcal F_M(U)$ agree on each $U_i$, then their $x$-coordinates agree for every $x\in U\cap Z$ because some $U_i$ contains $x$.  Hence the sections are equal.

Conversely, suppose $s_i\in\mathcal F_M(U_i)$ are compatible.  For each $x\in U\cap Z$, choose $i$ with $x\in U_i$ and define the $x$-coordinate of the glued section to be the $x$-coordinate of $s_i$.  Compatibility on overlaps makes this independent of the choice of $i$.  This produces the unique glued section on $U$.
\end{proof}

\begin{theorem}\label{thm:sheafification}
The canonical inclusion of presheaves
\[
\eta_D:\cP_D\longrightarrow \cV_D
\]
exhibits $\cV_D$ as the sheafification of $\cP_D$.
\end{theorem}

\begin{proof}
By Lemma~\ref{lem:product-sheaf}, $\cV_D$ is a sheaf.  It remains to prove the universal property.

Let $\mathcal F$ be any sheaf and let $\phi:\cP_D\to\mathcal F$ be a morphism of presheaves.  We construct a unique sheaf morphism $\widetilde\phi:\cV_D\to\mathcal F$ with $\widetilde\phi\circ\eta_D=\phi$.

Fix an open set $U\subset X$ and a section
\[
s=(s_x)_{x\in U\cap Z}\in\cV_D(U)=\prod_{x\in U\cap Z}A^{m_x}.
\]
For each $x\in U\cap Z$, choose an open neighborhood $U_x\subset U$ with $U_x\cap Z=\{x\}$.  Let $e_x\in\cP_D(U_x)$ be the section whose $x$-coordinate is $s_x$ and whose other coordinates are basepoints.  Set
\[
t_x:=\phi_{U_x}(e_x)\in\mathcal F(U_x).
\]
On an overlap $U_x\cap U_y$ with $x\neq y$, both $e_x$ and $e_y$ restrict to the unique basepoint section because $(U_x\cap U_y)\cap Z=\varnothing$.  Hence $t_x$ and $t_y$ agree on the overlap.  The family $(t_x)_{x\in U\cap Z}$ is therefore compatible, and together with the unique section on $U\setminus Z$ it glues to a section of $\mathcal F(U)$, which we define to be $\widetilde\phi_U(s)$.

This construction is functorial in $U$, so it defines a sheaf morphism $\widetilde\phi$.  By construction, if $s\in\cP_D(U)$ has finite support, the glued section is obtained from the same finitely many local pieces that define $\phi_U(s)$, so $\widetilde\phi_U(\eta_D(s))=\phi_U(s)$.  Thus $\widetilde\phi\circ\eta_D=\phi$.

Finally, uniqueness holds because every section of $\cV_D(U)$ is determined by its restrictions to the open cover
\[
U=(U\setminus Z)\cup\bigcup_{x\in U\cap Z}U_x,
\]
and on each $U_x$ the value of any factorization through a sheaf must agree with $\phi$ on the single-support section carrying the $x$-coordinate.
\end{proof}

\begin{proposition}\label{prop:shadow-stalks}
The sheaves $\cV_D$ and $\cA_D$ are well defined.  Their stalks at $x\in Z$ are
\[
(\cV_D)_x\cong A^{m_x},\qquad (\cA_D)_x\cong A,
\]
and their stalks away from $Z$ are one-point stalks.
\end{proposition}

\begin{proof}
This is Lemma~\ref{lem:product-sheaf}.  Since $Z$ is discrete, for each $x\in Z$ there is an open neighborhood $U$ with $U\cap Z=\{x\}$.  For every smaller neighborhood $V\subset U$ of $x$ one has
\[
\cV_D(V)=A^{m_x},\qquad \cA_D(V)=A,
\]
with identity restrictions, so the direct limits defining the stalks are $A^{m_x}$ and $A$.  If $x\notin Z$, choose $U$ with $U\cap Z=\varnothing$; then all sections on smaller neighborhoods are singleton products.
\end{proof}

\begin{proposition}\label{prop:boolean-cube}
For every integer $m\ge1$, the free rank-$m$ $A$-semimodule $A^m$ is canonically isomorphic to the Boolean cube
\[
B_m:=\cP([m]),\qquad [m]:=\{1,\dots,m\},
\]
by the rule
\[
(a_1,\dots,a_m)\longmapsto \{j:a_j=0\}.
\]
Under this identification, semimodule addition is union of subsets.
\end{proposition}

\begin{proof}
Because $A\cong\BB$, coordinatewise addition on $A^m$ is Boolean OR.  Therefore
\[
\{j:(a_j\oplus b_j)=0\}=\{j:a_j=0\text{ or }b_j=0\},
\]
which is exactly the union of the corresponding subsets.  Bijectivity is immediate.
\end{proof}

\begin{proposition}\label{prop:retract-shadow}
For each $m\ge1$, define
\[
\Sigma_m:A^m\to A,
\qquad
\Sigma_m(a_1,\dots,a_m):=a_1\oplus\cdots\oplus a_m,
\]
and
\[
\Delta_m:A\to A^m,
\qquad
\Delta_m(a):=(a,\dots,a).
\]
Then $\Sigma_m$ and $\Delta_m$ are $A$-semimodule maps and
\[
\Sigma_m\circ\Delta_m=\operatorname{id}_A.
\]
Consequently there are canonical sheaf morphisms
\[
\Sigma_D:\cV_D\to\cA_D,
\qquad
\Delta_D:\cA_D\to\cV_D,
\qquad
\Sigma_D\circ\Delta_D=\operatorname{id}_{\cA_D}.
\]
\end{proposition}

\begin{proof}
If $a=\taup$, then $\Sigma_m(\Delta_m(a))=\taup$.  If $a=0$, then idempotence in $A$ gives $\Sigma_m(\Delta_m(a))=0$.  The global statement is obtained coordinatewise over the support of the divisor.
\end{proof}

\begin{theorem}\label{thm:reduced-vs-multiplicity}
The following are equivalent.
\begin{enumerate}[label=(\roman*)]
\item $\cV_D\cong\cA_D$ in sheaves of $A$-semimodules.
\item $m_x=1$ for every $x\in Z$.
\end{enumerate}
When these conditions fail, $\cA_D$ is nevertheless a canonical split retract of $\cV_D$.
\end{theorem}

\begin{proof}
If every $m_x=1$, then the definitions give $\cV_D=\cA_D$.

Conversely, suppose $\cV_D\cong\cA_D$.  Then the stalks are isomorphic at every $x\in Z$, so
\[
A^{m_x}\cong A
\]
in $A$-semimodules.  Any such isomorphism is in particular a bijection of underlying sets, hence
\[
2^{m_x}=|A^{m_x}|=|A|=2.
\]
Thus $m_x=1$ for all $x$.  The split-retract statement is Proposition~\ref{prop:retract-shadow}.
\end{proof}

\begin{remark}
Because $\Spec(\BB)$ has one point, the Boolean divisor shadow is naturally a sheaf of semiring objects on a zero-dimensional support.  Through the embedding of semirings into blueprints \cite{Lorscheid}, it also admits a blueprint-theoretic interpretation.  We do not use that formalism below.
\end{remark}

\section{Classical Artin thickenings and the chain-cube comparison theorem}

\subsection{Local Artin rings}

\begin{definition}
For $x\in X$ and $m\ge1$, let $\mathfrak m_x\subset\cO_{X,x}$ be the maximal ideal.  Define the local Artin rings
\[
Q_{x,m}:=\cO_{X,x}/\mathfrak m_x^m,
\qquad
K_{x,m}:=\cO_{X,x}/\mathfrak m_x^{m+1}.
\]
For the divisor $D=\sum_{x\in Z}m_x[x]$, set
\[
Q_x:=Q_{x,m_x},\qquad K_x:=K_{x,m_x}.
\]
\end{definition}

\begin{lemma}\label{lem:dvr-local}
Let $x\in X$.  Then there exists a holomorphic coordinate $t$ centered at $x$ such that
\[
\cO_{X,x}\cong\CC\{t\},\qquad \mathfrak m_x=(t).
\]
Consequently
\[
Q_{x,m}\cong\CC\{t\}/(t^m)\cong\CC[t]/(t^m),
\qquad
K_{x,m}\cong\CC\{t\}/(t^{m+1})\cong\CC[t]/(t^{m+1}).
\]
\end{lemma}

\begin{proof}
The existence of a local holomorphic coordinate is part of the definition of a Riemann surface \cite[Chapter~1]{Lee}.  In such a coordinate, the germ ring is the convergent power-series ring $\CC\{t\}$ and the maximal ideal is generated by $t$.  Modulo $(t^m)$ every convergent power series is represented by its degree $<m$ truncation, which gives the polynomial description.
\end{proof}

\begin{proposition}\label{prop:ideal-chain-local}
For every $m\ge1$, the ideals of $Q_{x,m}$ are exactly
\[
I_k:=(t^k)/(t^m),\qquad k=0,1,\dots,m,
\]
where $t$ is any centered local coordinate.  In particular,
\[
Q_{x,m}=I_0\supset I_1\supset\cdots\supset I_m=0
\]
is a chain, independent of the chosen coordinate.
\end{proposition}

\begin{proof}
By Lemma~\ref{lem:dvr-local}, $Q_{x,m}\cong\CC\{t\}/(t^m)$.  Ideals of the quotient correspond to ideals of $\CC\{t\}$ containing $(t^m)$.  Since $\CC\{t\}$ is a discrete valuation ring, every ideal is $(t^k)$ for a unique $k\ge0$.  Those containing $(t^m)$ are precisely $(t^k)$ with $0\le k\le m$, giving the displayed chain.  Because the ideal lattice of a ring is intrinsic, the description is coordinate-independent.
\end{proof}

\begin{definition}
For $m\ge1$, define the terminal-segment chain
\[
C_m:=\{T_r:1\le r\le m+1\}\subset \cP([m]),
\]
where
\[
T_r:=\{r,r+1,\dots,m\}\quad (1\le r\le m),
\qquad
T_{m+1}:=\varnothing.
\]
We also write $B_m:=\cP([m])$ for the full Boolean cube.
\end{definition}

\begin{corollary}\label{cor:ideal-chain-identification}
For every $x\in X$ and $m\ge1$, the ideal lattice $\Idl(Q_{x,m})$ is canonically isomorphic to the chain $C_m$ via
\[
I_k=(t^k)/(t^m)\longmapsto T_{k+1}.
\]
\end{corollary}

\begin{proof}
Proposition~\ref{prop:ideal-chain-local} gives all ideals and their inclusion relations.  The displayed bijection preserves order because
\[
I_k\subseteq I_\ell\iff k\ge\ell
\iff T_{k+1}\subseteq T_{\ell+1}.
\]
\end{proof}

\subsection{The chain-cube comparison}

\begin{definition}
Define the inclusion
\[
i_m:C_m\into B_m
\]
by viewing each terminal segment as a subset of $[m]$.  Define monotone maps
\[
L_m:B_m\to C_m,
\qquad
L_m(I):=\begin{cases}
\varnothing,& I=\varnothing,\\
\{\min I,\min I+1,\dots,m\},& I\neq\varnothing,
\end{cases}
\]
and
\[
R_m:B_m\to C_m,
\qquad
R_m(I):=\bigcup\{T\in C_m:T\subseteq I\}.
\]
\end{definition}

\begin{lemma}\label{lem:LRtails}
For every $I\in B_m$, the sets $L_m(I)$ and $R_m(I)$ lie in $C_m$.  Moreover, $L_m(I)$ is the least terminal segment containing $I$, and $R_m(I)$ is the greatest terminal segment contained in $I$.
\end{lemma}

\begin{proof}
If $I=\varnothing$, then $L_m(I)=\varnothing=T_{m+1}\in C_m$.  If $I\neq\varnothing$, then $L_m(I)=T_{\min I}\in C_m$ and by construction it is the least terminal segment containing $I$.

For $R_m(I)$, note that the union of terminal segments contained in $I$ is again a terminal segment: if $T_r,T_s\subseteq I$, then $T_{\max\{r,s\}}\subseteq I$ and $T_r\cup T_s=T_{\min\{r,s\}}$.  Hence the displayed union belongs to $C_m$ and is plainly the greatest terminal segment contained in $I$.
\end{proof}

\begin{theorem}[Chain-cube comparison]\label{thm:chain-cube}
For every $m\ge1$ one has
\[
L_m\dashv i_m\dashv R_m
\]
in the category of finite posets.  Explicitly, for $I\in B_m$ and $T\in C_m$,
\[
L_m(I)\subseteq T \iff I\subseteq i_m(T),
\]
and
\[
i_m(T)\subseteq I \iff T\subseteq R_m(I).
\]
Hence $C_m$ is both reflective and coreflective in $B_m$.
\end{theorem}

\begin{proof}
By Lemma~\ref{lem:LRtails}, $L_m(I)$ is the least terminal segment containing $I$.  Therefore
\[
L_m(I)\subseteq T\iff I\subseteq T=i_m(T),
\]
which proves $L_m\dashv i_m$.

Similarly, $R_m(I)$ is the greatest terminal segment contained in $I$.  Therefore
\[
i_m(T)=T\subseteq I\iff T\subseteq R_m(I),
\]
which proves $i_m\dashv R_m$.
\end{proof}

\begin{corollary}\label{cor:not-iso-chain-cube}
If $m\ge2$, then $C_m$ and $B_m$ are not isomorphic in finite posets.  Nevertheless they are related by the biadjoint comparison of Theorem~\ref{thm:chain-cube}.
\end{corollary}

\begin{proof}
The chain $C_m$ has width $1$.  The Boolean cube $B_m$ has width at least $2$ because $\{1\}$ and $\{2\}$ are incomparable.  Poset isomorphisms preserve width, so no isomorphism exists when $m\ge2$.
\end{proof}

\subsection{Sheafified comparison}

\begin{definition}
Define sheaves of posets on $X$ by
\[
\cI_D(U):=\prod_{x\in U\cap Z}\Idl(Q_x),
\qquad
\cC_D(U):=\prod_{x\in U\cap Z}C_{m_x},
\qquad
\cB_D(U):=\prod_{x\in U\cap Z}B_{m_x},
\]
with coordinatewise order and restriction maps given by projection.
\end{definition}

\begin{proposition}\label{prop:sheaves-posets}
The assignments $\cI_D$, $\cC_D$, and $\cB_D$ are sheaves of finite posets.  There is a canonical isomorphism
\[
\cI_D\cong\cC_D.
\]
Moreover, under the identification of Proposition~\ref{prop:boolean-cube}, the ordered sheaf underlying the $A$-semimodule sheaf $\cV_D$ is canonically isomorphic to $\cB_D$.
\end{proposition}

\begin{proof}
This is Lemma~\ref{lem:product-sheaf} applied to families of finite posets, together with Corollary~\ref{cor:ideal-chain-identification} and Proposition~\ref{prop:boolean-cube} on each stalk.
\end{proof}

\begin{theorem}\label{thm:sheafified-chain-cube}
There are canonical sheaf morphisms of posets
\[
L_D:\cB_D\to\cC_D,
\qquad
i_D:\cC_D\to\cB_D,
\qquad
R_D:\cB_D\to\cC_D,
\]
obtained coordinatewise from $L_{m_x}$, $i_{m_x}$, and $R_{m_x}$.  Stalkwise one has
\[
(L_D)_x\dashv (i_D)_x\dashv (R_D)_x.
\]
Thus the classical jet-ideal chain sheaf is simultaneously a reflective and a coreflective subobject of the Boolean multiplicity-cube sheaf in sheaves of finite posets.
\end{theorem}

\begin{proof}
Coordinatewise application of Theorem~\ref{thm:chain-cube} commutes with restriction maps because all restrictions are coordinate projections.  Hence the maps are sheaf morphisms.  The stalkwise adjunctions are exactly Theorem~\ref{thm:chain-cube} at multiplicity $m_x$.
\end{proof}

\section{Derivative-enhanced jets and coefficient support}

\subsection{Coordinate-free local jet data}

\begin{definition}
Let $f\in\cO(X)$ be a nonzero holomorphic function with zero divisor
\[
\operatorname{div}(f)=D=\sum_{x\in Z}m_x[x].
\]
For each $x\in Z$, write $[f]_x$ for the class of $f$ in $K_x=\cO_{X,x}/\mathfrak m_x^{m_x+1}$.  Define the sheaves
\[
\cQ_D(U):=\prod_{x\in U\cap Z}Q_x,
\qquad
\cK_D(U):=\prod_{x\in U\cap Z}K_x,
\]
and the distinguished section
\[
[f]_D(U):=([f]_x)_{x\in U\cap Z}\in\cK_D(U).
\]
\end{definition}

\begin{proposition}\label{prop:derivative-enhanced-local}
For each $x\in Z$, the class $[f]_x$ is nonzero and lies in the minimal nonzero ideal
\[
\mathfrak m_x^{m_x}/\mathfrak m_x^{m_x+1}\subset K_x.
\]
Indeed, the principal ideal generated by $[f]_x$ is exactly this minimal nonzero ideal.
\end{proposition}

\begin{proof}
By definition of multiplicity, $\ord_x(f)=m_x$.  Hence in the discrete valuation ring $\cO_{X,x}$ one may write
\[
f=u_x\pi_x^{m_x}
\]
with $u_x$ a unit and $\pi_x$ a uniformizer.  Modulo $\mathfrak m_x^{m_x+1}$ this gives
\[
[f]_x=[u_x]\,[\pi_x]^{m_x},
\]
which is nonzero because $u_x$ remains invertible and $\pi_x^{m_x}\notin\mathfrak m_x^{m_x+1}$.  It lies in $\mathfrak m_x^{m_x}/\mathfrak m_x^{m_x+1}$ by construction.  Since that quotient is one-dimensional over the residue field, every nonzero element generates it, so the ideal generated by $[f]_x$ is exactly the minimal nonzero ideal.
\end{proof}

\begin{corollary}\label{cor:coordinate-derivative}
Let $x\in Z$ and choose a centered local coordinate $t$ at $x$.  Then in $K_x\cong\CC[t]/(t^{m_x+1})$ one has
\[
[f]_x=\lambda_x t^{m_x}
\qquad\text{with}\qquad
\lambda_x\in\CC^\times.
\]
If $X=\CC$ with global coordinate $s$, then
\[
\lambda_x=\frac{f^{(m_x)}(x)}{m_x!}.
\]
\end{corollary}

\begin{proof}
In the chosen coordinate, $f(t)=\lambda_x t^{m_x}+O(t^{m_x+1})$ with $\lambda_x\neq0$ because $\ord_x(f)=m_x$.  Passing to the quotient by $(t^{m_x+1})$ gives the first formula.  On $\CC$, this is Taylor's theorem in the coordinate $s$.
\end{proof}

\subsection{Coefficient-support maps on $\CC$}

From now until the end of this subsection we specialize to $X=\CC$ with global coordinate $s$.

\begin{definition}
For $x\in\CC$ and $m\ge1$, define the derivative-enhanced stalk
\[
E_{x,m}:=(s-x)\CC[s]/(s-x)^{m+1}.
\]
Its distinguished basepoint is the zero class.  Under the basis
\[
(s-x),(s-x)^2,\dots,(s-x)^m,
\]
each element of $E_{x,m}$ can be written uniquely as
\[
\sum_{r=1}^m c_r(s-x)^r.
\]
Define a pointed-set map
\[
\sigma_{x,m}:E_{x,m}\to A^m
\]
by
\[
\sigma_{x,m}\!\left(\sum_{r=1}^m c_r(s-x)^r\right)=(a_1,\dots,a_m),
\qquad
 a_r:=\begin{cases}0,& c_r\neq0,\\ \taup,& c_r=0.\end{cases}
\]
Define also
\[
\iota_{x,m}:A^m\to E_{x,m}
\]
by
\[
\iota_{x,m}(a_1,\dots,a_m):=\sum_{\{r:a_r=0\}}(s-x)^r.
\]
\end{definition}

\begin{proposition}\label{prop:pointed-retract}
The maps $\sigma_{x,m}$ and $\iota_{x,m}$ are pointed-set morphisms and satisfy
\[
\sigma_{x,m}\circ\iota_{x,m}=\operatorname{id}_{A^m}.
\]
Consequently, if
\[
\cE_D(U):=\prod_{x\in U\cap Z}E_{x,m_x}
\]
on $\CC$, then the coordinatewise maps
\[
\sigma_D:\cE_D\to\cV_D,
\qquad
\iota_D:\cV_D\to\cE_D
\]
exhibit the underlying pointed-set sheaf of $\cV_D$ as a retract of $\cE_D$.
\end{proposition}

\begin{proof}
Both maps preserve the chosen basepoints by construction.  If $a=(a_1,\dots,a_m)\in A^m$, then the coefficient of $(s-x)^r$ in $\iota_{x,m}(a)$ is nonzero exactly when $a_r=0$.  Applying $\sigma_{x,m}$ therefore recovers $a$.  The sheaf statement is obtained coordinatewise and uses Lemma~\ref{lem:product-sheaf}.
\end{proof}

\begin{proposition}\label{prop:not-cmon-iso}
If $m\ge1$, then $E_{x,m}$ and $A^m$ are not isomorphic as commutative monoids under addition.
\end{proposition}

\begin{proof}
The additive monoid of $A^m$ is idempotent because addition is coordinatewise Boolean OR:
\[
u+u=u\qquad (u\in A^m).
\]
The additive monoid of $E_{x,m}$ is not idempotent: for example,
\[
(s-x)+(s-x)=2(s-x)\neq (s-x)
\]
in characteristic $0$.  Idempotence of the entire monoid is preserved by monoid isomorphism, so no isomorphism exists.
\end{proof}

\begin{corollary}\label{cor:zeta-support-slot}
Let $\rho$ be a zero of zeta of multiplicity $m_\rho$.  In the stalk
\[
E_\rho:=E_{\rho,m_\rho}
\]
one has
\[
[\zeta]_\rho=\lambda_\rho (s-\rho)^{m_\rho},
\qquad
\lambda_\rho:=\frac{\zeta^{(m_\rho)}(\rho)}{m_\rho!}\neq0.
\]
Under the support map $\sigma_{\rho,m_\rho}$, this class is sent to the element of $A^{m_\rho}$ whose only nonbasepoint coordinate is the terminal one.  Equivalently, it picks the atom $\{m_\rho\}\subset [m_\rho]$ of the terminal-segment chain.
\end{corollary}

\begin{proof}
The displayed formula is Corollary~\ref{cor:coordinate-derivative} applied to $f=\zeta$.  Since only the coefficient of $(s-\rho)^{m_\rho}$ is nonzero, the support vector has a single nonbasepoint entry in the terminal slot.
\end{proof}

\section{Compactification and tail objects at infinity}

Let
\[
j:\CC\into\PP^1(\CC)
\]
be the standard open immersion.

\subsection{Pointed-set tails}

\begin{definition}
Let $Z\subset\CC$ be closed and discrete, and let $(M_x)_{x\in Z}$ be pointed sets with basepoints $0_x$.  Define the sheaf
\[
\mathcal F_M(U):=\prod_{x\in U\cap Z}M_x
\qquad (U\subset\CC\text{ open}).
\]
For an open set $U\subset\PP^1(\CC)$ containing $\infty$, write $U^\circ:=U\cap\CC$.  Define
\[
\Tail((M_x)_{x\in Z}):=\prod_{x\in Z}M_x\Big/\sim,
\]
where
\[
(a_x)\sim(b_x)\iff a_x=b_x\text{ for all but finitely many }x\in Z.
\]
\end{definition}

\begin{theorem}\label{thm:stalks-at-infty}
With the notation above, the canonical morphism
\[
j_!\mathcal F_M\longrightarrow j_*\mathcal F_M
\]
has the following stalks.
\begin{enumerate}[label=(\roman*)]
\item If $x\in Z$, then
\[
(j_!\mathcal F_M)_x\cong (j_*\mathcal F_M)_x\cong M_x.
\]
\item If $y\in\CC\setminus Z$, then
\[
(j_!\mathcal F_M)_y\cong (j_*\mathcal F_M)_y\cong *.
\]
\item At $\infty$ one has
\[
(j_!\mathcal F_M)_\infty\cong *,
\qquad
(j_*\mathcal F_M)_\infty\cong \Tail((M_x)_{x\in Z}).
\]
\end{enumerate}
Consequently, if the tail pointed set is nontrivial, then
\[
j_!\mathcal F_M\not\cong j_*\mathcal F_M
\]
in sheaves of pointed sets on $\PP^1(\CC)$.
\end{theorem}

\begin{proof}
The first two assertions are immediate from sufficiently small neighborhoods in $\CC$: around a point of $Z$ the support is a singleton, while around a point of $\CC\setminus Z$ the support is empty.

For the stalk of $j_!\mathcal F_M$ at $\infty$, note that a section of $j_!\mathcal F_M(U)$ on a neighborhood $U$ of $\infty$ is, by definition, eventually equal to the basepoint outside a finite subset of $U^\circ\cap Z$.  Shrinking $U$ so that those finitely many points disappear shows that every germ becomes the basepoint germ.

For the stalk of $j_*\mathcal F_M$ at $\infty$, a germ is represented by a section on some neighborhood $U$ of $\infty$, hence by a family on the cofinite subset $U^\circ\cap Z\subset Z$.  Extend that family by basepoints on the omitted finite subset of $Z$.  This identifies the germ with a class in the quotient modulo finite modification.  Two such families define the same germ exactly when they agree on a smaller neighborhood of $\infty$, equivalently when they differ only at finitely many indices.  This gives the stated identification.

If the stalks at $\infty$ are not isomorphic, the sheaves are not isomorphic.
\end{proof}

\subsection{Vector-space tails}

\begin{definition}
Let $(V_x)_{x\in Z}$ be a family of nonzero complex vector spaces.  Define the sheaf
\[
\mathcal F_V(U):=\prod_{x\in U\cap Z}V_x
\qquad (U\subset\CC\text{ open})
\]
and the tail vector space
\[
\Tail_{\mathrm{vec}}((V_x)_{x\in Z}):=\prod_{x\in Z}V_x\Big/\bigoplus_{x\in Z}V_x.
\]
\end{definition}

\begin{proposition}\label{prop:vector-tail-exact}
There is a short exact sequence of sheaves of complex vector spaces on $\PP^1(\CC)$,
\[
0\longrightarrow j_!\mathcal F_V
\longrightarrow j_*\mathcal F_V
\longrightarrow i_{\infty *}\Tail_{\mathrm{vec}}((V_x)_{x\in Z})
\longrightarrow0.
\]
If $Z$ is infinite, then the tail vector space is nonzero.  Moreover,
\[
H^q\bigl(\PP^1(\CC),j_!\mathcal F_V\bigr)=0
\qquad (q\ge1).
\]
\end{proposition}

\begin{proof}
On an open set $U\subset\PP^1(\CC)$ not containing $\infty$, the sheaves $j_!\mathcal F_V$ and $j_*\mathcal F_V$ coincide.  If $\infty\in U$, then
\[
(j_*\mathcal F_V)(U)=\prod_{x\in U^\circ\cap Z}V_x,
\qquad
(j_!\mathcal F_V)(U)=\bigoplus_{x\in U^\circ\cap Z}V_x,
\]
because sections of $j_!$ are precisely those with finite support near $\infty$.  Since $U^\circ\cap Z$ is cofinite in $Z$, the quotient is canonically identified with
\[
\prod_{x\in Z}V_x\Big/\bigoplus_{x\in Z}V_x.
\]
This proves exactness and identifies the quotient with a skyscraper at $\infty$.

If $Z$ is infinite and each $V_x\neq0$, any family $(v_x)$ with all $v_x\neq0$ defines a nonzero class in the quotient, so the tail vector space is nonzero.

The sheaves $j_*\mathcal F_V$ and $i_{\infty*}\Tail_{\mathrm{vec}}((V_x))$ are flabby: restriction maps on products are coordinate projections, hence surjective, and the same is immediate for a skyscraper sheaf.  The long exact sequence in sheaf cohomology therefore gives
\[
H^q(\PP^1(\CC),j_!\mathcal F_V)=0\qquad(q\ge2),
\]
and an exact segment
\[
0\to \Gamma(j_!\mathcal F_V)\to \Gamma(j_*\mathcal F_V)\to \Tail_{\mathrm{vec}}((V_x))\to H^1(\PP^1(\CC),j_!\mathcal F_V)\to0.
\]
But
\[
\Gamma(j_!\mathcal F_V)=\bigoplus_{x\in Z}V_x,
\qquad
\Gamma(j_*\mathcal F_V)=\prod_{x\in Z}V_x,
\]
and the middle map is the quotient map defining the tail vector space, hence is surjective.  Therefore $H^1$ vanishes as well.
\end{proof}

\section{The zeta divisor, packet descent, and trivial-zero weights}

\subsection{The zeta divisor}

\begin{definition}
Let
\[
D_\zeta:=\sum_{\rho:\,\zeta(\rho)=0}m_\rho[\rho]
\]
be the effective divisor of zeta zeros on $\CC$.  Set
\[
\cP_\zeta:=\cP_{D_\zeta},\quad
\cV_\zeta:=\cV_{D_\zeta},\quad
\cA_\zeta:=\cA_{D_\zeta},\quad
\cI_\zeta:=\cI_{D_\zeta},\quad
\cQ_\zeta:=\cQ_{D_\zeta},\quad
\cK_\zeta:=\cK_{D_\zeta}.
\]
\end{definition}

\begin{corollary}\label{cor:zeta-shadow-sheaves}
At every zero $\rho$ of zeta,
\[
(\cV_\zeta)_\rho\cong A^{m_\rho},
\qquad
(\cA_\zeta)_\rho\cong A,
\qquad
\Fix_S(F_1(\rho))=A.
\]
The reduced zeta shadow is a canonical split retract of the multiplicity shadow.
\end{corollary}

\begin{proof}
Combine Proposition~\ref{prop:shadow-stalks}, Proposition~\ref{prop:retract-shadow}, and Corollary~\ref{cor:zeta-pointwise}.
\end{proof}

\subsection{Completed zeta and packet descent}

\begin{definition}
Define the completed zeta function
\[
\xi(s):=\frac12 s(s-1)\pi^{-s/2}\Gamma\!\left(\frac s2\right)\zeta(s).
\]
Let $D_\xi$ be its zero divisor.  Thus $D_\xi$ is the divisor of the nontrivial zeros of zeta.
\end{definition}

\begin{definition}
Let
\[
c(s):=\overline s,
\qquad
w(s):=1-s,
\qquad
W:=\{\mathrm{id},c,w,cw\}\cong C_2\times C_2.
\]
\end{definition}

\begin{lemma}\label{lem:xi-symmetry-multiplicity}
If $\rho$ is a zero of $\xi$, then $\overline\rho$, $1-\rho$, and $1-\overline\rho$ are also zeros of $\xi$, and all four points occur with the same multiplicity.
\end{lemma}

\begin{proof}
The functional equation $\xi(1-s)=\xi(s)$ gives the reflection symmetry.  The Dirichlet series for $\zeta$ has real coefficients on $\Re(s)>1$, so $\overline{\zeta(\overline s)}=\zeta(s)$ there, and analytic continuation yields the same relation meromorphically on all of $\CC$.  The elementary prefactor in $\xi$ is also conjugation-compatible, so $\overline{\xi(\overline s)}=\xi(s)$.  Therefore zeros are preserved by conjugation.  Multiplicities are preserved because both symmetries are local analytic automorphisms with nonzero derivative in absolute value.
\end{proof}

\begin{definition}
Let $Y_\xi:=|D_\xi|/W$ be the orbit space of nontrivial-zero packets, equipped with the quotient topology.  Since $|D_\xi|$ is discrete, so is $Y_\xi$.  For an orbit $O\in Y_\xi$, write $m_O$ for the common multiplicity of its points.

Define the \emph{Boolean packet sheaf}
\[
\mathcal P_\xi^{\cV}(U):=\prod_{O\in U}A^{m_O}
\]
and the \emph{chain packet sheaf}
\[
\mathcal P_\xi^{\cC}(U):=\prod_{O\in U}C_{m_O}
\]
on open sets $U\subset Y_\xi$.
\end{definition}

\begin{theorem}\label{thm:packet-descent}
Let $q:|D_\xi|\to Y_\xi$ be the quotient map.  The group $W$ acts naturally on the sheaves $\cV_{D_\xi}$ and $\cC_{D_\xi}$.  For each orbit $O\in Y_\xi$ of cardinality $r_O$, the stalk of $q_*\cV_{D_\xi}$ at $O$ is canonically
\[
(q_*\cV_{D_\xi})_O\cong (A^{m_O})^{r_O},
\]
and its $W$-fixed subsemimodule is the diagonal copy of $A^{m_O}$.  Hence
\[
(q_*\cV_{D_\xi})^W\cong \mathcal P_\xi^{\cV}.
\]
Likewise,
\[
(q_*\cC_{D_\xi})_O\cong (C_{m_O})^{r_O},
\qquad
(q_*\cC_{D_\xi})^W\cong \mathcal P_\xi^{\cC}.
\]
Thus both the Boolean multiplicity shadow and the classical jet-ideal chain descend to packet sheaves on the orbit space.
\end{theorem}

\begin{proof}
Because $Y_\xi$ is discrete, the stalk of $q_*\cV_{D_\xi}$ at an orbit $O$ is the product of the stalks over the points in that orbit.  Lemma~\ref{lem:xi-symmetry-multiplicity} gives the common multiplicity $m_O$, so
\[
(q_*\cV_{D_\xi})_O\cong \prod_{\rho\in O}A^{m_O}\cong (A^{m_O})^{r_O}.
\]
The group $W$ acts by permuting the $r_O$ factors transitively.  A tuple in $(A^{m_O})^{r_O}$ is invariant under this action if and only if all coordinates are equal.  Thus the fixed part is the diagonal copy of $A^{m_O}$.  Since the orbit space is discrete, these diagonal stalks assemble into the sheaf $\mathcal P_\xi^{\cV}$.

The chain case is identical, replacing $A^{m_O}$ by $C_{m_O}$.
\end{proof}

\begin{remark}
No hypothesis such as RH is used here.  If RH holds, every nontrivial-zero packet has cardinality $2$; off the critical line, the packet size is $4$.
\end{remark}

\subsection{Trivial zeros}

\begin{proposition}\label{prop:trivial-simple}
For every integer $n\ge1$, the trivial zero $-2n$ of zeta is simple.
\end{proposition}

\begin{proof}
The functional equation reads
\[
\zeta(s)=2^s\pi^{s-1}\sin\!\left(\frac{\pi s}{2}\right)\Gamma(1-s)\zeta(1-s).
\]
At $s=-2n$, the factors $2^s$, $\pi^{s-1}$, and $\Gamma(1-s)=\Gamma(1+2n)$ are holomorphic and nonzero, while $\zeta(1-s)=\zeta(2n+1)\neq0$.  The only vanishing factor is the sine term, and it has a simple zero because
\[
\frac{d}{ds}\sin\!\left(\frac{\pi s}{2}\right)\Big|_{s=-2n}=\frac\pi2\cos(-n\pi)=\frac\pi2(-1)^n\neq0.
\]
\end{proof}

\begin{corollary}\label{cor:trivial-stalk}
On the trivial-zero sector, the Boolean multiplicity stalk is just $A$:
\[
(\cV_\zeta)_{-2n}\cong A\qquad(n\ge1).
\]
\end{corollary}

\begin{proof}
This is Proposition~\ref{prop:shadow-stalks} together with Proposition~\ref{prop:trivial-simple}.
\end{proof}

\begin{proposition}\label{prop:trivial-derivative}
For every integer $n\ge1$,
\[
\zeta'(-2n)=(-1)^n\frac{(2n)!}{2(2\pi)^{2n}}\,\zeta(2n+1).
\]
Hence
\[
|\zeta'(-2n)|=\frac{(2n)!}{2(2\pi)^{2n}}\,\zeta(2n+1).
\]
\end{proposition}

\begin{proof}
Set
\[
g_n(s):=2^s\pi^{s-1}\Gamma(1-s)\zeta(1-s).
\]
Then $\zeta(s)=g_n(s)\sin(\pi s/2)$ and $g_n$ is holomorphic at $s=-2n$ with
\[
g_n(-2n)=2^{-2n}\pi^{-2n-1}(2n)!\,\zeta(2n+1).
\]
Differentiating and using that the sine factor vanishes at $-2n$ gives
\[
\zeta'(-2n)=g_n(-2n)\cdot \frac\pi2\cos(-n\pi)
=2^{-2n}\pi^{-2n-1}(2n)!\,\zeta(2n+1)\cdot\frac\pi2(-1)^n,
\]
which simplifies to the stated formula.  Since $\zeta(2n+1)>0$, the absolute-value formula follows.
\end{proof}

\begin{definition}
For each trivial zero $-2n$, define its \emph{canonical Boolean weight}
\[
\omega_{-2n}:A\to\RR_{\ge0}
\]
by
\[
\omega_{-2n}(\taup):=0,
\qquad
\omega_{-2n}(0):=|\zeta'(-2n)|.
\]
\end{definition}

\begin{corollary}\label{cor:trivial-weight}
The unique nonbasepoint of the trivial Boolean stalk carries the exact analytic weight
\[
\omega_{-2n}(0)=\frac{(2n)!}{2(2\pi)^{2n}}\,\zeta(2n+1).
\]
\end{corollary}

\begin{proof}
This is immediate from the definition and Proposition~\ref{prop:trivial-derivative}.
\end{proof}

\subsection{Compactified zeta tails}

\begin{definition}
Define the zeta tail objects
\[
T_\zeta^{\cV}:=\Tail\bigl((A^{m_\rho})_{\rho:\,\zeta(\rho)=0}\bigr),
\qquad
T_\zeta^{\cA}:=\Tail\bigl((A)_{\rho:\,\zeta(\rho)=0}\bigr),
\]
and the jet-vector-space tails
\[
T_\zeta^{\cQ}:=\Tail_{\mathrm{vec}}\bigl((Q_\rho)_{\rho:\,\zeta(\rho)=0}\bigr),
\qquad
T_\zeta^{\cK}:=\Tail_{\mathrm{vec}}\bigl((K_\rho)_{\rho:\,\zeta(\rho)=0}\bigr).
\]
\end{definition}

\begin{theorem}\label{thm:zeta-tail}
The compactified Boolean and jet shadows of zeta satisfy:
\begin{enumerate}[label=(\roman*)]
\item the stalks at $\infty$ are
\[
(j_*\cV_\zeta)_\infty\cong T_\zeta^{\cV},
\qquad
(j_*\cA_\zeta)_\infty\cong T_\zeta^{\cA},
\]
\[
(j_*\cQ_\zeta)_\infty\cong T_\zeta^{\cQ},
\qquad
(j_*\cK_\zeta)_\infty\cong T_\zeta^{\cK};
\]
\item all four tail objects are nonzero;
\item for the jet-vector-space shadows there are short exact sequences
\[
0\to j_!\cQ_\zeta\to j_*\cQ_\zeta\to i_{\infty*}T_\zeta^{\cQ}\to0,
\]
\[
0\to j_!\cK_\zeta\to j_*\cK_\zeta\to i_{\infty*}T_\zeta^{\cK}\to0,
\]
and therefore
\[
H^q\bigl(\PP^1(\CC),j_!\cQ_\zeta\bigr)=H^q\bigl(\PP^1(\CC),j_!\cK_\zeta\bigr)=0
\qquad(q\ge1).
\]
\end{enumerate}
\end{theorem}

\begin{proof}
Part~(i) is Theorem~\ref{thm:stalks-at-infty} for Boolean shadows and Proposition~\ref{prop:vector-tail-exact} for jet shadows.  Since zeta has infinitely many trivial zeros, the index set is infinite; the class of the everywhere-nonbasepoint family on the trivial-zero sector shows that the Boolean tails are nonzero, and any everywhere-nonzero vector in the product gives nonzero jet tails.  This proves~(ii).  Part~(iii) is Proposition~\ref{prop:vector-tail-exact}.
\end{proof}

\section{Cyclotomic phase coordinates and orthogonal walks}\label{sec:tobley}

In this section we record two exact phase-coordinate constructions extracted from unpublished notes of Jacolm Tobley \cite{Tobley}.  They are independent of the preceding sheaf-theoretic arguments, but they apply naturally to the normalized leading coefficients of simple local zeta jets.

\begin{proposition}[Orthogonal-walk expansion]\label{prop:orthogonal-walk}
Let $\alpha_1,\dots,\alpha_m\in\RR$, and let $e_r(\alpha_1,\dots,\alpha_m)$ be the elementary symmetric polynomials.  Then
\[
\prod_{k=1}^m (1+i\alpha_k)=\sum_{r=0}^m i^r e_r(\alpha_1,\dots,\alpha_m).
\]
Consequently, if
\[
S_r:=\sum_{j=0}^r i^j e_j(\alpha_1,\dots,\alpha_m),
\qquad 0\le r\le m,
\]
then the successive differences $S_r-S_{r-1}$ are alternately horizontal and vertical segments of lengths $e_r(\alpha_1,\dots,\alpha_m)$.  In particular,
\[
\prod_{k=1}^m \frac{1+i\alpha_k}{\sqrt{1+\alpha_k^2}}
=e^{i\sum_{k=1}^m\arctan(\alpha_k)}.
\]
\end{proposition}

\begin{proof}
Expanding the product gives
\[
\prod_{k=1}^m(1+i\alpha_k)=\sum_{I\subseteq [m]} i^{|I|}\prod_{j\in I}\alpha_j.
\]
Grouping terms by the cardinality $r=|I|$ yields the first formula because
\[
e_r(\alpha_1,\dots,\alpha_m)=\sum_{|I|=r}\prod_{j\in I}\alpha_j.
\]
The geometric statement follows from the fact that multiplication by $i^r$ rotates the positive real axis through successive quarter-turns.  Finally,
\[
\frac{1+i\alpha_k}{\sqrt{1+\alpha_k^2}}
=\cos(\arctan\alpha_k)+i\sin(\arctan\alpha_k)
=e^{i\arctan\alpha_k},
\]
and multiplying over $k$ proves the normalized identity.
\end{proof}

\begin{theorem}[Cyclotomic affine-factor chart]\label{thm:cyclotomic-chart}
Let $\zeta_5=e^{2\pi i/5}$ and set $\phi:=2\pi/5$.  For $t>0$, define
\[
L(t):=1+t\zeta_5.
\]
Then the function
\[
\Theta:(0,\infty)\to(0,\phi),\qquad \Theta(t):=\Arg L(t),
\]
is a real-analytic strictly increasing bijection.  Its inverse is
\[
t(\theta)=\frac{\sin\theta}{\sin(\phi-\theta)}
\qquad (0<\theta<\phi).
\]
Equivalently,
\[
\frac{1+t(\theta)\zeta_5}{|1+t(\theta)\zeta_5|}=e^{i\theta}
\qquad (0<\theta<2\pi/5).
\]
\end{theorem}

\begin{proof}
Write $z(t):=1+t e^{i\phi}$.  Since $\phi\in(0,\pi)$, one has $\Im z(t)=t\sin\phi>0$, so $\Arg z(t)$ indeed lies in $(0,\phi)$.  Differentiating the argument gives
\[
\Theta'(t)=\Im\!\left(\frac{z'(t)}{z(t)}\right)
=\Im\!\left(\frac{e^{i\phi}}{1+t e^{i\phi}}\right)
=\frac{\sin\phi}{|1+t e^{i\phi}|^2}>0.
\]
Thus $\Theta$ is strictly increasing and real-analytic.  Moreover,
\[
\lim_{t\to0^+}\Theta(t)=0,
\qquad
\lim_{t\to\infty}\Theta(t)=\phi,
\]
so $\Theta$ is a bijection from $(0,\infty)$ onto $(0,\phi)$.

Now let $0<\theta<\phi$ and seek $t>0$ such that
\[
\Arg(1+t e^{i\phi})=\theta.
\]
This means that $1+t e^{i\phi}$ lies on the ray $\RR_{>0}e^{i\theta}$, so there exists $r>0$ with
\[
1+t e^{i\phi}=r e^{i\theta}.
\]
Taking imaginary parts yields
\[
t\sin\phi=r\sin\theta,
\]
and taking real parts yields
\[
1+t\cos\phi=r\cos\theta.
\]
Eliminating $r$ gives
\[
\frac{t\sin\phi}{\sin\theta}=\frac{1+t\cos\phi}{\cos\theta},
\]
whence
\[
t\sin(\phi-\theta)=\sin\theta.
\]
Therefore
\[
t=\frac{\sin\theta}{\sin(\phi-\theta)}.
\]
The final normalized identity is just a restatement of the defining equation for $t(\theta)$.
\end{proof}

\begin{corollary}\label{cor:simple-zero-phase-chart}
Let $\rho$ be a simple zero of zeta, and define its normalized leading coefficient
\[
u_\rho:=\frac{\zeta'(\rho)}{|\zeta'(\rho)|}\in S^1.
\]
If $\Arg(u_\rho)\in(0,2\pi/5)$, then there exists a unique $t_\rho>0$ such that
\[
u_\rho=\frac{1+t_\rho\zeta_5}{|1+t_\rho\zeta_5|}.
\]
Thus the phase of the simple local zeta jet is carried by a unique normalized affine factor in the CM field $\QQ(\zeta_5)$ on that arc.
\end{corollary}

\begin{proof}
This is an immediate application of Theorem~\ref{thm:cyclotomic-chart} with $\theta=\Arg(u_\rho)$.
\end{proof}

\begin{remark}
Theorem~\ref{thm:cyclotomic-chart} is purely archimedean.  It does not assert that zeta derivatives are cyclotomic numbers.  Rather, it provides an exact phase chart whenever a chosen normalized local coefficient lies in the indicated arc.
\end{remark}

\section{Final theorem}

\begin{theorem}
Let
\[
S=G(\ZZ)=\ZZ\sqcup\{\taup\},\qquad A=\{0,\taup\}\cong\BB,
\qquad F_1(s)=\zeta(s)-1.
\]
Then the zeta zero divisor admits the following canonical structures.
\begin{enumerate}[label=(\roman*)]
\item For every zero $\rho$ of zeta,
\[
F_1(\rho)=-1,
\qquad
\Fix_S(F_1(\rho))=A.
\]
\item For every effective divisor $D$ on a Riemann surface, the finite-support presheaf $\cP_D$ sheafifies to the Boolean multiplicity shadow $\cV_D$, and the reduced shadow $\cA_D$ is a canonical split retract of $\cV_D$.  One has $\cV_D\cong\cA_D$ if and only if every multiplicity is $1$.
\item The local Artin thickening $Q_{x,m_x}=\cO_{X,x}/\mathfrak m_x^{m_x}$ has ideal lattice $C_{m_x}$, the Boolean shadow stalk is the cube $B_{m_x}$, and the precise comparison is the biadjunction
\[
L_{m_x}\dashv i_{m_x}\dashv R_{m_x}.
\]
\item For a holomorphic function $f$ with divisor $D$, the local jet class in $K_{x,m_x}=\cO_{X,x}/\mathfrak m_x^{m_x+1}$ generates the minimal nonzero ideal.  On $\CC$, the derivative-enhanced pointed-set sheaf retracts onto the Boolean multiplicity shadow, while no additive-monoid isomorphism exists between the local jet stalk and the Boolean cube.
\item After compactification to $\PP^1(\CC)$, the difference between direct image and extension by zero is measured by an explicit tail object at $\infty$.  For zeta these tails are nonzero.
\item For the completed zeta function $\xi$, the nontrivial-zero shadows descend to packet sheaves on the Klein-four orbit space.
\item On the trivial-zero sector, the Boolean stalk is simple and carries the exact weight
\[
|\zeta'(-2n)|=\frac{(2n)!}{2(2\pi)^{2n}}\,\zeta(2n+1).
\]
\item The normalized leading coefficients of simple local zeta jets admit exact cyclotomic phase coordinates on the arc $(0,2\pi/5)$, and the product of normalized affine factors is governed by the elementary-symmetric expansion of Proposition~\ref{prop:orthogonal-walk}.
\end{enumerate}
\end{theorem}

\begin{proof}
The individual assertions are proved in Corollary~\ref{cor:zeta-pointwise}, Theorem~\ref{thm:sheafification}, Theorem~\ref{thm:reduced-vs-multiplicity}, Theorem~\ref{thm:chain-cube}, Proposition~\ref{prop:derivative-enhanced-local}, Proposition~\ref{prop:pointed-retract}, Proposition~\ref{prop:not-cmon-iso}, Theorem~\ref{thm:stalks-at-infty}, Theorem~\ref{thm:packet-descent}, Proposition~\ref{prop:trivial-derivative}, Theorem~\ref{thm:zeta-tail}, Proposition~\ref{prop:orthogonal-walk}, and Theorem~\ref{thm:cyclotomic-chart}.
\end{proof}

\begin{thebibliography}{99}

\bibitem{DrosteKuichVogler}
M.~Droste, W.~Kuich, and H.~Vogler (eds.),
\emph{Handbook of Weighted Automata},
Springer, Berlin, 2009.

\bibitem{GaubertKatz}
S.~Gaubert and R.~Katz,
\emph{Rational semimodules over the max-plus semiring and geometric approach of discrete event systems},
arXiv:math/0208014, 2002.

\bibitem{GrothendieckEGAII}
A.~Grothendieck,
\emph{\'El\'ements de g\'eom\'etrie alg\'ebrique. II. \`Etude globale \`el\'ementaire de quelques classes de morphismes},
Publ. Math. IH\'ES \textbf{8} (1961), 5--222.

\bibitem{Gunawardena}
J.~Gunawardena (ed.),
\emph{Idempotency},
Cambridge University Press, Cambridge, 1998.

\bibitem{JarraLorscheidVital}
M.~Jarra, O.~Lorscheid, and E.~Vital,
\emph{Quiver matroids, matroid morphisms, quiver Grassmannians, their Euler characteristics and $\mathbb F_1$-points},
arXiv:2404.09255v2, 2026.

\bibitem{Lee}
J.~M.~Lee,
\emph{Introduction to Complex Manifolds},
American Mathematical Society, Providence, RI, 2024.

\bibitem{Lorscheid}
O.~Lorscheid,
\emph{The geometry of blueprints. Part I: Algebraic background and scheme theory},
Adv. Math. \textbf{229} (2012), no.~3, 1804--1846.

\bibitem{OrecchiaChiantini}
F.~Orecchia and L.~Chiantini (eds.),
\emph{Zero-Dimensional Schemes},
Walter de Gruyter, Berlin, 1994.

\bibitem{Tobley}
J.~Tobley,
\emph{Unpublished notes on cyclotomic affine phase factors, arctangent decompositions, and orthogonal-walk expansions},
May 2026.  Communicated in the project source transcripts supplied with the manuscript materials.

\bibitem{Vakil}
R.~Vakil,
\emph{The Rising Sea: Foundations of Algebraic Geometry},
2024 draft.

\end{thebibliography}

\end{document}
